%%%%%%%%%%%%%%%%%%%%%%%%%%%%%%%%%
%
%       EXERCÍCIO
%
%%%%%%%%%%%%%%%%%%%%%%%%%%%%%%%%%

\ifdefstring{\atividade}{prova}{%
    \ifdefstring{\modo}{objetivo}{%
        \renewcommand{\valorquestao}{\ValorQObj\ ponto}
    }{%
        \renewcommand{\valorquestao}{\ValorQDisc\ pontos}
    }%
}%

\begin{exercicioBanco}[\valorquestao]
Um pedestre caminha em uma avenida retilínea. Sua distância, em quilômetros, até uma farmácia localizada no quilômetro \(1\) é dada por
\[
|x-1|,
\]
e sua distância até um hospital localizado no quilômetro \(6\) é dada por
\[
|x-6|.
\]

Sabe-se que a soma dessas distâncias é igual a \(5,5\) quilômetros. Determine as posições possíveis do pedestre.


% Define as alternativas
\newcommand{\alternativas}{%
    \begin{center}
        \begin{tabularx}{\textwidth}{XXXXX}
            (a) \(\{0,7\}\). &
            (b) \(\left\{\dfrac{3}{4},\dfrac{25}{4}\right\}\). &
            (c) \(\left\{\dfrac{11}{4},\dfrac{11}{2}\right\}\). &
            (d) \(\{0,11\}\). &
            (e) \(\{7,25\}\).
        \end{tabularx}
    \end{center}
}

% Define a resposta correta
\newcommand{\resposta}{B}

% Lógica condicional para exibição
\ifdefstring{\atividade}{lista}{%
    \alternativas
    \vspace{0.5em}
    
    \noindent\textbf{Resposta:} letra \textbf{\resposta}.
}{%
    \ifdefstring{\modo}{objetivo}{%
        \alternativas
    }{%
        % modo = discursiva → não mostra alternativas
    }
}
\end{exercicioBanco}

